\uextra{Apéndice}{Especificación del Protocolo de Comunicación con micro:bit}
\label{apx:microbit-protocol}

Este documento define un protocolo para enviar subflujos desde NodeLand al micro:bit. El objetivo del protocolo es que el host pueda desplegar un subflujo, ejecutarlo con argumentos, y recibir resultado/logs/eventos, de forma robusta en enlaces con MTU pequeña (BLE) y posibles fragmentaciones.

\subsection*{1) Transporte}

\begin{itemize}
    \item \textbf{Medio}: USB Serial (Web Serial) o BLE UART (microbit-web-bluetooth).
    \item \textbf{Codificación}: UTF-8 / ASCII.
    \item \textbf{Framing}: 1 mensaje equivale a 1 línea terminada en \texttt{\textbackslash{}n}.
    \item En Bluetooth, el host envía cada línea en chunks de 18 caracteres para evitar saturación del buffer.
\end{itemize}

\subsection*{2) Formato de frame (línea)}

Todos los mensajes del protocolo empiezan por el prefijo \texttt{DFP1|}.

\begin{verbatim}
DFP1|<TYPE>|<SEQ>|<PAYLOAD>
\end{verbatim}

\begin{itemize}
    \item \texttt{TYPE}: string en mayúsculas (sin \texttt{|}).
    \item \texttt{SEQ}: entero decimal \texttt{0..65535} para ACK.
    \item \texttt{PAYLOAD}: lista de campos clave-valor (\texttt{k=v}) separados por \texttt{;}.
          \begin{itemize}
              \item \texttt{k}: \texttt{[A-Z0-9\_]+}
              \item \texttt{v}: por defecto ASCII simple (sin \texttt{;}), y para binario usar \textbf{HEX} (base16).
          \end{itemize}
\end{itemize}

Notas de parsing:
\begin{itemize}
    \item El payload no debe contener \texttt{|}.
    \item Si un valor requiere \texttt{;} o \texttt{=}, debe enviarse en HEX.
\end{itemize}

Ejemplo:
\begin{verbatim}
DFP1|PING|1|
DFP1|HELLO_RES|2|DEV=MBIT-123;FW=0.1.0;
CAPS=ALU,DISPLAY;MAX_PROG=2048;MAX_CHUNK=32
\end{verbatim}

\subsection*{3) ACK, reintentos y chunking}

\subsubsection*{ACK}

Para mensajes críticos (BEGIN/CHUNK/COMMIT, RUN, STOP) el receptor responde:

\begin{itemize}
    \item OK:\\
          \texttt{DFP1|ACK|<SEQ>|OK=1}
    \item Error:\\
          \texttt{DFP1|ACK|<SEQ>|OK=0;ERR=<CODE>;MSG\_HEX=<hex\_utf8>}
\end{itemize}

\texttt{CODE} recomendado: \texttt{BAD\_FORMAT}, \texttt{UNSUPPORTED}, \texttt{CRC\_FAIL}, \texttt{OUT\_OF\_MEMORY}, \texttt{UNKNOWN\_PROGRAM}, \texttt{BUSY}.

\subsubsection*{Chunking (envío de bytecode)}

El flujo para el envío de programas sigue los siguientes pasos:

\begin{enumerate}
    \item \texttt{B} (BEGIN, metadatos)
    \item \texttt{C} (CHUNK, repetido)
    \item \texttt{M} (COMMIT, verificación final)
\end{enumerate}

Sintaxis de cada tipo:
\begin{itemize}
    \item \texttt{B} (BEGIN): \texttt{DFP1|B|0|L=<len\_base36>;C=<chunk\_bytes>}
    \item \texttt{C} (CHUNK): \texttt{DFP1|C|0|X=O<off\_base36>H<hex>}
    \item \texttt{M} (COMMIT): \texttt{DFP1|M|0|R=<crc32\_base36>}
    \item \texttt{A} (ARGS chunk): \texttt{DFP1|A|0|I=<run\_id>;X=<hex>}
    \item \texttt{R} (RUN): \texttt{DFP1|R|<SEQ>|I=<run\_id>}
\end{itemize}

Parámetros de envío recomendados:
\begin{itemize}
    \item Deploy chunk: \texttt{1 byte} por frame (\texttt{C}).
    \item Args chunk: \texttt{1 byte} por frame (\texttt{A}, es decir \texttt{X} de 2 chars hex).
\end{itemize}

Timeouts/reintentos utilizados:
\begin{itemize}
    \item \texttt{B/C/A}: timeout 3000ms, 1 reintento
    \item \texttt{M}: timeout 5000ms, 1 reintento
    \item \texttt{R}: timeout 90000ms
\end{itemize}

\subsection*{4) Control de ejecución}

\subsubsection*{RUN}
\begin{enumerate}
    \item Uno o más \texttt{A}:
          \begin{verbatim}
DFP1|A|0|I=<run_id>;X=<hex>
\end{verbatim}
    \item Luego \texttt{R} para ejecutar:
          \begin{verbatim}
DFP1|R|<SEQ>|I=<run_id>
\end{verbatim}
\end{enumerate}

Donde \texttt{I} es el \texttt{run\_id}. El firmware reconstruye \texttt{ARGS\_HEX} concatenando los \texttt{X} recibidos para ese \texttt{run\_id}.

\subsubsection*{STOP}

Si el host quiere detener una ejecución en curso:
\begin{verbatim}
DFP1|STOP|21|I=<run_id>
\end{verbatim}

\subsubsection*{RESULT}

Resultado:
\begin{verbatim}
DFP1|RESULT|0|I=<run_id>;VAL_HEX=<hex>
\end{verbatim}


\texttt{SEQ} puede ser \texttt{0} en mensajes espontáneos desde micro:bit (no ACK).

\subsection*{5) Handshake y capacidades}

\begin{itemize}
    \item Host $\rightarrow$ micro:bit:
          \begin{verbatim}
    DFP1|HELLO_REQ|1|
    \end{verbatim}

    \item micro:bit $\rightarrow$ host:
          \begin{verbatim}
    DFP1|HELLO_RES|1|DEV=<id>;FW=<semver>;
    CAPS=<csv>;MAX_PROG=<bytes>;MAX_CHUNK=<bytes>
    \end{verbatim}
\end{itemize}

\texttt{CAPS} es una lista de capacidades soportadas por el firmware, que indica qué comandos/operaciones están disponibles. Se utilizaron las siguientes capacidades para el prototipo inicial:
\begin{itemize}
    \item \texttt{ALU} (operaciones numéricas)
    \item \texttt{STR} (strings; concat, replace, subtract, slice)
    \item \texttt{DISPLAY} (matriz LED)
    \item \texttt{SOUND}
    \item \texttt{PINS}
    \item \texttt{BUTTONS}
    \item \texttt{SENSORS}
    \item \texttt{LIST}
    \item \texttt{LOGIC}
\end{itemize}

El host puede compilar/validar el subflujo según \texttt{CAPS}.

\subsection*{6) Formato binario (bytecode)}

Cada instrucción: \texttt{[op:u8][a:u8][b:u8][c:u8]}.

Opcodes:
\begin{itemize}
    \item \texttt{0x01 PUSH\_CONST a=constIndex}
    \item \texttt{0x02 LOAD\_ARG a=argIndex}
    \item \texttt{0x10 ADD} / \texttt{0x11 SUB} / \texttt{0x12 MUL} / \texttt{0x13 DIV} / \texttt{0x14 MOD} / \texttt{0x15 POW}
    \item \texttt{0x16 FLOOR} / \texttt{0x17 CEIL} / \texttt{0x18 ROUND}
    \item \texttt{0x20 CONCAT}
    \item \texttt{0x21 REPLACE\_TEXT} / \texttt{0x22 SUBTRACT\_TEXT} / \texttt{0x23 SLICE\_TEXT}
    \item \texttt{0x30 DISPLAY\_TEXT} (pop string)
    \item \texttt{0x31 DISPLAY\_NUMBER} (pop number)
    \item \texttt{0x32 PLAY\_TONE} (pop ms, pop freq)
    \item \texttt{0x34 READ\_TEMP}
    \item \texttt{0x50 LIST\_COUNT} / \texttt{0x51 LIST\_SORT\_ASC} / \texttt{0x52 LIST\_SORT\_DESC}
    \item \texttt{0x53 LIST\_ACCUMULATE} / \texttt{0x54 LIST\_ZIP} / \texttt{0x55 LIST\_UNZIP\_A} / \texttt{0x56 LIST\_UNZIP\_B}
    \item \texttt{0x60 CMP\_EQ} / \texttt{0x61 CMP\_NEQ} / \texttt{0x62 CMP\_LT} / \texttt{0x63 CMP\_LTE} / \texttt{0x64 CMP\_GT} / \texttt{0x65 CMP\_GTE}
    \item \texttt{0x70 LOGIC\_AND} / \texttt{0x71 LOGIC\_NAND} / \texttt{0x72 LOGIC\_OR} / \texttt{0x73 LOGIC\_NOR} / \texttt{0x74 LOGIC\_XOR} / \texttt{0x75 LOGIC\_NOT}
    \item \texttt{0x40 RETURN} (pop value)
\end{itemize}

\subsection*{7) Argumentos y valores (VAL\_HEX / ARGS\_HEX)}
Para \texttt{ARGS\_HEX} y \texttt{VAL\_HEX} se usa una codificación binaria simple:
\begin{itemize}
    \item \texttt{tag:u8} + payload
          \begin{itemize}
              \item \texttt{0=number\_f32} + 4 bytes
              \item \texttt{1=string} + \texttt{u16 len} + bytes
              \item \texttt{2=bool} + \texttt{u8}
              \item \texttt{3=null}
              \item \texttt{4=list} + \texttt{u8 len} + valores anidados
          \end{itemize}
\end{itemize}

\texttt{ARGS\_HEX} contiene:
\begin{itemize}
    \item \texttt{argc:u8} seguido de \texttt{argc} valores.
\end{itemize}

\texttt{VAL\_HEX} contiene un solo valor.